%!TEX TS-program = xelatex
%!TEX encoding = UTF-8 Unicode
% Awesome CV LaTeX Template for CV/Resume
% 
% Template license:
% CC BY-SA 4.0 (https://creativecommons.org/licenses/by-sa/4.0/)
%
% Adaptation by
% Lara J. Martin (laramar@seas.upenn.edu)
%
% Template from 
% https://github.com/Denubis/Academic-LaTeX-CV-BibTeX-CSV
% Brian Ballsun-Stanton (brian.ballsun-stanton@mq.edu.au)
%
% Originally taken from
% https://github.com/posquit0/Awesome-CV
% Claud D. Park <posquit0.bj@gmail.com>
% http://www.posquit0.com
%-------------------------------------------------------------------------------
% CONFIGURATIONS
%-------------------------------------------------------------------------------
% A4 paper size by default, use 'letterpaper' for US letter
\documentclass[11pt, a4paper]{awesome-cv}


\usepackage{datatool}
\usepackage{xparse}
\usepackage{eso-pic,graphicx}
\usepackage{pstricks}
\usepackage{hyphenat}

    
% https://tex.stackexchange.com/a/376174

% Configure page margins with geometry
\geometry{left=1.4cm, top=.8cm, right=1.4cm, bottom=1.8cm, footskip=.5cm}

% Specify the location of the included fonts
\fontdir[fonts/]

% Color for highlights
% Awesome Colors: awesome-emerald, awesome-skyblue, awesome-red, awesome-pink, awesome-orange
%                 awesome-nephritis, awesome-concrete, awesome-darknight
%\colorlet{awesome}{awesome-red}
% Uncomment if you would like to specify your own color
% \definecolor{awesome}{HTML}{CA63A8}

% Colors for text
% Uncomment if you would like to specify your own color
% \definecolor{darktext}{HTML}{414141}
% \definecolor{text}{HTML}{333333}
% \definecolor{graytext}{HTML}{5D5D5D}
% \definecolor{lighttext}{HTML}{999999}

% Set false if you don't want to highlight section with awesome color
\setbool{acvSectionColorHighlight}{false}

% If you would like to change the social information separator from a pipe (|) to something else
\renewcommand{\acvHeaderSocialSep}{\quad\textbar\quad}

\newcommand{\myname}{\textbf{L. J. Martin}}
\newcommand{\mynamedagger}{\textbf{L. J. Martin$^\dagger$}}

% Following entry pattern: https://zety.com/blog/academic-cv-example


%-------------------------------------------------------------------------------
%	PERSONAL INFORMATION
%	Comment any of the lines below if they are not required
%-------------------------------------------------------------------------------
% Available options: circle|rectangle,edge/noedge,left/right
% \photo{./examples/profile.png}

\addbibresource{bibliography.bib}

%%%%%%%%%%%%%%%%%%%%%%
% Sort Bibliography
% https://tex.stackexchange.com/questions/484108/biblatex-sort-by-year-and-name-both-descending
\DeclareSortingTemplate{ydndt}{
  \sort{
    \field{presort}
  }
  \sort[final]{
    \field{sortkey}
  }
  \sort[direction=descending]{
    \field{sortyear}
    \field{year}
    \literal{9999}
  }
  \sort[direction=descending]{
    \field{sortmonth}
    \field{month}
    \literal{12}
  }
  \sort{
    \field{sorttitle}
    \field{title}
  }
}
%%%%%%%%%%%%%%%%%%%%%%
% Sort citation numbers
% https://tex.stackexchange.com/questions/273313/biblatex-sort-cites-via-number-split-bibliography
% https://tex.stackexchange.com/questions/310684/how-to-get-sorted-and-compressed-citations-when-using-split-bibliography-lists-a
\DeclareSourcemap{
  \maps[datatype=bibtex]{
    \map{
     \step[fieldsource=keywords, match=thesis, final]
     \step[fieldset=presort, fieldvalue = {A}]
    }
    \map{
      \step[fieldsource=keywords, match=conference, final]
      \step[fieldset=presort, fieldvalue = {B}]
    }
    \map{
      \step[fieldsource=keywords, match=workshop, final]
      \step[fieldset=presort, fieldvalue = {C}]
    }
    \map{
      \step[fieldsource=keywords, match=other, final]
      \step[fieldset=presort, fieldvalue = {D}]
    }
    \map{
      \step[fieldsource=keywords, match=preprint, final]
      \step[fieldset=presort, fieldvalue = {E}]
    }
    \map{
      \step[fieldsource=keywords, match=submitted, final]
      \step[fieldset=presort, fieldvalue = {F}]
    }
  }
}
%%%%%%%%%%%%%%%%%%%%%%




% Available options: circle|rectangle,edge/noedge,left/right
% \photo{./examples/profile.png}
\name{Dr.}{Lara J.}{Martin}
\position{Assistant Professor, CSEE \bullet~University of Maryland, Baltimore County (UMBC)}
\address{Information Technology and Engineering (ITE) Building, Room 216}
%\position{CIFellow Postdoctoral Researcher \bullet~University of Pennsylvania}
%\address{200 South 33rd Street, Room 297, Philadelphia, PA 19104}

%\mobile{(+1) 609-975-9799}
\homepage{laramartin.net}
\email{laramar@umbc.edu}
\github{lara-martin}
\linkedin{lara-j-martin}
% \gitlab{gitlab-id}
% \stackoverflow{SO-id}{SO-name}
\twitter{@LangTechLara}
\mastodon{https://sigmoid.social/@laramar}
% \skype{skype-id}
% \reddit{reddit-id}
% \medium{madium-id}
\googlescholar{YjiWURYAAAAJ}{Google Scholar}
%% \firstname and \lastname will be used
% \googlescholar{googlescholar-id}{}
% \extrainfo{extra informations}


%-------------------------------------------------------------------------------
\begin{document}

\AddToShipoutPictureBG*{%
  \AtPageUpperLeft{%
    \raisebox{-\height}{%
      \includegraphics[width=\paperwidth,height=0.175\paperwidth]{sidebar.png}%
    }%
  }
}

% Print the header with above personal information
% Give optional argument to change alignment(C: center, L: left, R: right)
\makecvheader[C]



% Print the footer with 3 arguments(<left>, <center>, <right>)
% Leave any of these blank if they are not needed
\makecvfooter
  {Last updated: \today}
  {Dr. Lara J. Martin~~~·~~~Curriculum Vitae}
  {\thepage}


%-------------------------------------------------------------------------------
%	CV/RESUME CONTENT
%	Each section is imported separately, open each file in turn to modify content
%-------------------------------------------------------------------------------
\vspace{1cm}
{\color{black}\fontsize{12pt}{1em}\faIcon{search}} \cvsection{Research Interests}

%\entrytitlestyle{Human-Centered AI}

%generally through neurosymbolic methods

\headerquotestyle{Human-Centered Artificial Intelligence, Computational Creativity, Automated Story Generation \& Understanding, Neurosymbolic Methods, Augmentative \& Alternative Communication (AAC), Speech Processing, Affective Computing, Conversational Agents, Cognitive Systems}


%Data scientist, educator, sysadmin, and data security specialist with over nine years designing and delivering technical solutions for academic and student research projects at the Macquarie University Faculty of Arts and UNSW Australia. Chief Investigator in category one and two grants, contracts, and prizes across the humanities, social sciences, and security studies totalling over AUD\$3,248,624. Lead Investigator on a big data investigation of violent extremism using tens of millions of posts on social media deploying computational data collection and analysis techniques using Machine Learning and Natural Language Processing. Presented findings for this project at the highest levels of state and federal governments, leading to international collaborations with NATO and the United States Institute for Peace. Technical Director, Product Owner, and DevOps/SysAdmin for a field-data collection project, delivering 64+ field data collection modules since 2013. Educator, trainer, and mentor staff and postgraduate students across the Faculty of Arts conducting digitally-enabled research. Trainer and Instructor for The Carpentries, leading over eleven workshops in data management, reproducible research techniques, and pedagogical techniques. Convener of undergraduate and postgraduate units in Information Technology and Digital Humanities topics. Supervised two MRes students applying and assessing Machine Learning techniques in sociology and ancient history. Strong industry ties with Google, Ubisoft, and Twitter -- bringing Macquarie University on as the first university on Google Arts and Culture, invited to be an academic advisor for the Fabricius Workbench, and facilitating big data social media collection on Twitter. Provision of technical advice, ethics and security advice, data management plans, database designs, and solution-focused support for academic and student research projects across Australian and International academic institutions. An internationally respected researcher in the Digital Humanities with sixteen peer-reviewed publications, commissioned reports, newspaper articles, book chapters, and conference presentations.




\vspace{0.5cm}
%-------------------------------------------------------------------------------
%	SECTION TITLE
%-------------------------------------------------------------------------------
{\color{black}\fontsize{12pt}{1em}\faIcon{graduation-cap}} \cvsection{ Education}


%-------------------------------------------------------------------------------
%	CONTENT
%-------------------------------------------------------------------------------
\begin{cventries}

%---------------------------------------------------------
  \cventry
    {Ph.D. in Human-Centered Computing} % Degree
    {Georgia Institute of Technology} % Institution
    {Aug. 2015 - May 2021} %– Dec. 2020} % Date(s)
    {Atlanta, GA} % Location
    {
      \begin{cvitems} % Description(s) bullet points
       % \item{School of Interactive Computing, College of Computing}
        \item {Advisor: Dr. Mark O. Riedl}
        \item {Thesis: \href{https://smartech.gatech.edu/handle/1853/64643}{Neurosymbolic Automated Story Generation}}
        \item{Teaching Certification: \href{https://ctl.gatech.edu/content/tech-teaching-0}{Tech to Teaching} (Summer 2018)}
      \end{cvitems}
    }



 \cventry
    {M.S. in Language Technologies} % Degree
    {Carnegie Mellon University} % Institution
    {Aug. 2013 – Aug. 2015} % Date(s)
    {Pittsburgh, PA} % Location
    {      
      \begin{cvitems} % Description(s) bullet points
        \item {Advisor: Dr. Alan W Black}
      \end{cvitems}
    }


  \cventry
    {B.S. in Computer Science \& Linguistics (double major)} % Degree
    {Rutgers University --- New Brunswick} % Institution
    {Sep. 2009 – May 2013} % Date(s)
    {Piscataway, NJ} % Location
    {      
      \begin{cvitems} % Description(s) bullet points
      \item {Advisor: Dr. Matthew Stone}
      \end{cvitems}
    }%
%---------------------------------------------------------
\end{cventries}

\vspace{0.5cm}
%-------------------------------------------------------------------------------
%	SECTION TITLE
%-------------------------------------------------------------------------------
{\color{black}\fontsize{12pt}{1em}\faIcon{building}} \cvsection{Research Experience}


%-------------------------------------------------------------------------------
%	CONTENT
%-------------------------------------------------------------------------------
\begin{cventries}


%---------------------------------------------------------
  \cventry
    {University of Maryland, Baltimore County -- Computer Science and Electrical Engineering (CSEE) Department} % Organization
    {Assistant Professor} % Job title
    {Aug 2023 – Present} % Date(s)
    {Baltimore, MD} % Location
    {
      % \begin{cvitems} % Description(s) of tasks/responsibilities
      % \end{cvitems}
    }

%---------------------------------------------------------
  \cventry
    {University of Pennsylvania -- Computer and Information Science} % Organization
    {\href{https://cifellows2020.org/}{Computing Innovation Fellow (CIFellow)} Postdoctoral Researcher} % Job title
    {Jan 2021 – Aug 2023} % Date(s)
    {Philadelphia, PA} % Location
    {
      % \begin{cvitems} % Description(s) of tasks/responsibilities
      % \item {Identifying the story understanding capabilities of large language models (LLMs).}
      % \item {Developing a working AAC prototype given feedback from users.}
      % \item {Conducted semi-structured interviews with autistic adult users of augmentative and alternative communication (AAC).}
      % \end{cvitems}
    }
%---------------------------------------------------------
  \cventry
    {Georgia Institute of Technology -- School of Interactive Computing} % Organization
    {Graduate Research Assistant} % Job title
    {Aug 2015 – Dec 2020} % Date(s)
    {Atlanta, GA} % Location
    {
      % \begin{cvitems} % Description(s) of tasks/responsibilities
      %   \item {Created a complex end-to-end automated story generation pipeline.}
      % \end{cvitems}
    }

%---------------------------------------------------------
  \cventry
    {Amazon.com Inc. -- Alexa Smart Home Machine Learning} % Organization
    {Applied Scientist Intern} % Job title
    {May 2017 – Aug 2017} % Date(s)
    {Seattle, WA} % Location
    {
      % \begin{cvitems} % Description(s) of tasks/responsibilities
      %   \item {Identified potential research questions within Alexa Smart Home.}
      %   \item {Developed a system for identifying commands with an assumed context.}
      % \end{cvitems}
    }

%---------------------------------------------------------
  \cventry
   {Carnegie Mellon University -- Language Technologies Institute} % Organization
    {Graduate Research Assistant} % Job title
    {Sept 2013 – Aug 2015} % Date(s)
    {Pittsburgh, PA} % Location
    {
      % \begin{cvitems} % Description(s) of tasks/responsibilities
      %   \item {Created a zero-resource speech-to-speech translation system for the University of Pittsburgh Medical Center.}
      %   \item {Performed emotion recognition in noisy speech for event detection.}
      % \end{cvitems}
    }
%---------------------------------------------------------
  \cventry
   {University of Southern California -- Institute for Creative Technologies} % Organization
    {Intern} % Job title
    {May 2011 – Aug 2011} % Date(s)
    {Playa Vista, CA} % Location
    {
      % \begin{cvitems} % Description(s) of tasks/responsibilities
      %   \item {Wrote a chatbot for the Virtual Patient Project using Bruce Wilcox’s language Chatscript.}
      %   \item {Developed an authoring tool for the Chatscript language using Java.}
      %   \item {Designed and ran experiments comparing my Chatscript system to the project’s current chat system.}
      % \end{cvitems}
    }

%---------------------------------------------------------
\end{cventries}

\vspace{0.5cm}

\DTLloaddb[]{teaching}{cv/teaching.csv}

%\newpage

{\color{black}\fontsize{12pt}{1em}\faIcon{chalkboard-teacher}} \cvsection{Teaching}

\begin{cventries}
\DTLforeach{teaching}{%
%ProjectName    Date    Role    Location    Description1 Link  
\projectname=ProjectName,%
\date=Date,%
\role=Role,%
\location=Location,%
\descone=Description1,%
\link=Link%
}%
{%
    \DTLifnullorempty{\link}
    {
    \cventryNoNote
    {\role} % subhead
    {\projectname} % head
    {\date} % green
    {\location} % flushright grey
    }%
    {%
    \cventry
    {\role} % subhead
    {\href{\link}{\projectname}} % head
    {\date} % green
    {\location} % flushright grey
    {%
    \DTLifnullorempty{\descone}{}
    {
    \begin{cvitems}
    \item \descone%
    \end{cvitems}}
    }%
    }%
}
\end{cventries}

\vspace{.5cm}
{\color{black}\fontsize{12pt}{1em}\faIcon{box}} \cvsection{Projects}

% \DTLloaddb[]{projects}{cv/projects.csv}

% \cvsubsection{Summary}
% \def\projectaffiliated{0}
% \def\projectexternal{0}
% \def\projecttotal{0}

% \DTLforeach{projects}{%
% %GrantAffiliated	External
% \affiliated=GrantAffiliated,%
% \external=External}{%
% \DTLgadd{\projectaffiliated}{\affiliated}{\projectaffiliated}%
% \DTLgadd{\projectexternal}{\external}{\projectexternal}%
% \DTLgadd{\projecttotal}{1}{\projecttotal}%
% }
% \begin{cvhonors}
% \cvhonor{projects with external funding}{excluding FAIMS sub-projects}{2012 -- Present}{\projectaffiliated{}}
% \cvhonor{projects with external impact}{excluding FAIMS sub-projects}{2012 -- Present}{\projectexternal{}}
% \cvhonor{total projects}{excluding FAIMS sub-projects}{2012 -- Present}{\projecttotal{}}
% \cvhonor{FAIMS sub-projects}{modules, extensions, and custom exporters}{2014 -- Present}{64}
% \end{cvhonors}
% \cvsubsection{Project Details}
% \begin{cventries}
% \DTLforeach{projects}{%
% %ProjectName	StartYear	StopYear	Collaborator	Tools	Description1	Description2
% \projectname=ProjectName,%
% \startyear=StartYear,%
% \stopyear=StopYear,%
% \collaborator=Collaborator,%
% \tools=Tools,%
% \descone=Description1,%
% \desctwo=Description2}%
% {%
%  \cventry
%     {\tools} % subhead
%     {\projectname} % head
%     {\startyear{} -- \DTLifnullorempty{\stopyear}{Present}{\stopyear{}}} % green
%     {\collaborator} % flushright grey
%     {%
%     %\ifthenelse{\DTLcurrentindex=16}{\dtlbreak}{}%
%     \DTLifnullorempty
%         {\descone}
%         {}
%         {\begin{cvitems}
%         \item \descone%
%         \DTLifnullorempty
%             {\desctwo}
%             {}
%             {\item \desctwo}
%         \end{cvitems}
%         }
%     }
% }



% \end{cventries}

{\color{teal}\faIcon{icons} {\bf Human-AI Communication and Computational Creativity.}} {\bodyfontlight Collaboration between people and AI, primarily through storytelling.}

\begin{tabular*}{\textwidth}{@{\extracolsep{\fill}}p{.2cm} L{6cm} L{\textwidth - 10cm} @{\extracolsep{\fill}}R{3cm}}
    \setlength{\extrarowheight}{5pt}
    \faIcon{book-open} & Story understanding &  \cite{Li2022,Giorgi2023,Dong2023} & \entrylocationstyle{2021 -- Present}\\
    \faIcon{d-and-d} & D\&D as an AI challenge & \cite{Zhu2023AIIDE,martin2018dungeons,CallisonBurch2022,Papazov2022, CallisonBurchEMNLP,Zhu2023} & \entrylocationstyle{2018 -- Present}\\
    \faIcon{edit} & Story generation &  \cite{martin2016improvisational,martin2017improvisational,martin2017event,martin2018event,ammanabrolu2019guided,tambwekar2019controllable,ammanabrolu2020story,martin2021thesis,Alabdulkarim2021} & \entrylocationstyle{2015 -- Present}\\
    \faIcon{bezier-curve} & Schema co-creation & \cite{Zhang2023} & \entrylocationstyle{2023}\\
    \faIcon{comments} & Conversational agents & \cite{panagopoulouquakerbot}, See also: Amazon \& USC ICT internships & \entrylocationstyle{2011, 2017, 2022}\\
    %mobile 
    \faIcon{headphones} & Language learning & \cite{wolfeapplause} & \entrylocationstyle{2014}\\
    
\end{tabular*}

\vspace{.1cm}

{\color{teal}\faIcon{people-arrows} {\bf Human-Human Communication.}} {\bodyfontlight Analysis of human-human communication or computer-mediated communication.}



\begin{tabular*}{\textwidth}{@{\extracolsep{\fill}} p{.2cm} L{9cm} L{5.6cm} R{2.2cm}}
    \setlength{\extrarowheight}{5pt}
    \faIcon{tty} & Augmentative and Alternative Communication (AAC) & In progress. & \entrylocationstyle{2022 -- Present}\\
    \faIcon{reddit-alien} & Online communities  &\cite{moon2014identifying,Giorgi2023} & \entrylocationstyle{2014, 2023}\\
    \faIcon{language} & Translation & \cite{martin2015utterance}& \entrylocationstyle{2015}\\
    \faIcon{frown-open} & Emotion recognition \& affective computing &  \cite{martin2014methodology,cosentino2014,Yu2014,Yu2015}& \entrylocationstyle{2013 -- 2015}\\
    
    
\end{tabular*}


\vspace{.5cm}
%-------------------------------------------------------------------------------
%	SECTION TITLE
%-------------------------------------------------------------------------------
{\color{black}\fontsize{12pt}{1em}\faIcon{paper-plane}} \cvsection{Publications}
\vspace{0.5em}
\defbibnote{bothnote}{ {\footnotesize $^*$equal contribution, $^\dagger$presented}}
\defbibnote{presented}{ {\footnotesize $^\dagger$presented}}

\nocite{*}
\printbibliography[keyword={thesis}, title=PhD Dissertation]
\vspace{1em}
\printbibliography[prenote=bothnote, keyword={conference}, title=Conference Proceedings]
\vspace{1em}
\printbibliography[prenote=bothnote, keyword={workshop}, title=Refereed Workshop Papers]
%\vspace{1em}
\newpage
\printbibliography[keyword={other}, title=Other Publications]
\vspace{1em}
\printbibliography[keyword={preprint}, title=Preprints]
\printbibliography[keyword={submitted}, title=Submitted]



\vspace{.5cm}
{\color{black}\fontsize{12pt}{1em}\faIcon{file-invoice-dollar}} \cvsection{Grants and Awards}

\DTLloaddb[]{grants}{cv/grants.csv}

% \cvsubsection{Grant Summary}

% \def\granttotal{0}
% \def\grantnum{0}
% \DTLforeach{grants}{%
% \amount=Amount}{%
% \DTLgadd{\granttotal}{\amount}{\granttotal}%
% \DTLgadd{\grantnum}{1}{\grantnum}%
% }

% \begin{cvhonors}
% \cvhonor{total external competitive grant amount awarded}{}{2021 -- Present}{\$\granttotal{}}
% \cvhonor{grants awarded}{}{2021 -- Present}{\grantnum{}}
% \end{cvhonors}

\cvsubsection{Grants}

\begin{cventries}

\DTLforeach{grants}{%
%Grant	Code	Title	Amount	Start Year	End Year	Role	Location	Description	Primary CI	Authors
\grant=Organization,%
\code=Code,%
\granttitle=Title,%
\amount=Amount,%
\startyear=StartYear,%
\stopyear=EndYear,%
\role=Role,%
\location=Location,%
\desc=Description,%
\authors=Authors,%
\link=Link%
}{%
\cvgrant
    {\DTLifnullorempty{\code}{}{(\code)~}\DTLifnullorempty{\link}{\granttitle}{\href{\link}{\granttitle}}}%
    {\grant}%
    {\startyear{}\DTLifnullorempty{\stopyear}{}{ -- \stopyear{}}}%
    {\$\DTLgadd{\amount}{0}{\amount}\amount}%
    {\authors}%
    {\location}%
    {}}%
\end{cventries}




% \cvgrant
%     {\DTLifnullorempty{\code}{}{(\code)~}\granttitle{}}%
%     {\grant}%
%     {\startyear{} -- \stopyear{}}%
%     {AUD \DTLdecimaltocurrency{\income}{\grantincome}\grantincome}%
%     {\authors}%
%     {\location}%
%     {\desc{} Lead: \pci{}.}}% (\role{}) 

\cvsubsection{Awards}

\begin{cventries}
\cvgrant
{\href{https://gvu.gatech.edu/gvu-graduate-student-awards-program-2019}{Foley Scholar Award}}% title
{James D. Foley GVU Center Endowment}% grant agency
{2019}
{\$5,000}
{Lara J. Martin}
{Georgia Institute of Technology}{}
%{Every year, three PhD students across Georgia Tech's College of Computing are selected from a group of around 8 finalists to receive the prestigious Foley Scholarship in recognition for their efforts in innovative research.}




% \cvgrant
% {Travel Scholarship for the ACM Richard Tapia Celebration of Diversity in Computing}% title
% {Association for Computing Machinery}% grant agency
% {2019}
% {\$500}
% {Lara J. Martin}
% {San Diego, CA}
% {}

% \cvgrant
% {Travel Scholarship for the 2013 Grace Hopper Celebration of Women in Computing}% title
% {Lockheed Martin Corporation}% grant agency
% {2013}
% {\$450}
% {Lara J. Martin}
% {Minneapolis, MN}
% {}

\end{cventries}
\vspace{.5cm}
{\color{black}\fontsize{12pt}{1em}\faIcon{users}} \cvsection{Institutional Service}

\begin{cventries}
\cventry
    {Reviewer}
    {President's Undergraduate Research Awards (PURA)}
    {Summer 2019}
    {Georgia Institute of Technology}
    {}    
\cventry
    {Volunteer}
    {School of Interactive Computing’s Prospective Student Visit Week}
    {Spring '16, '17, '18}
    {Georgia Institute of Technology}
    {}
\cventry
    {Member}
    {School of Interactive Computing Faculty Hiring Committee}
    {Fall 2018}
    {Georgia Institute of Technology}
    {}
\cventry
    {Member}
    {Graduate Student Council}
    {Fall 2018-Spring 2019}
    {Georgia Institute of Technology}
    {}    
\cventry
    {Website Manager}
    {Human-Centered Computing Website}
    {Fall 2017-Spring 2019}
    {Georgia Institute of Technology}
    {}    
\cventry
    {Coordinator}
    {School of Interactive Computing's Prospective Student Visit Week}
    {Spring 2016}
    {Georgia Institute of Technology}
    {}
\end{cventries}



\vspace{.5cm}
%-------------------------------------------------------------------------------
%	SECTION TITLE
%-------------------------------------------------------------------------------

{\color{black}\fontsize{12pt}{1em}\faIcon{comments}} \cvsection{Professional Activities -- Conference Organization}

\DTLloaddb[]{organizing}{cv/organizing.csv}
\cvsubsection{Summary}
\def\organizer{0}
\def\chair{0}
\def\reviewer{0}
\def\attendee{0}
\def\totalconf{0}


\DTLforeach{organizing}{%
\status=GeneralRole}{%
\DTLifeq{\status}{Organizer}{%
\DTLgadd{\organizer}{1}{\organizer}%
}{}%
\DTLifeq{\status}{Action Editor}{%
\DTLgadd{\chair}{1}{\chair}%
}{}%
\DTLifeq{\status}{Chair}{%
\DTLgadd{\chair}{1}{\chair}%
}{}%
\DTLifeq{\status}{Reviewer}{%
\DTLgadd{\reviewer}{1}{\reviewer}%
}{}%
\DTLifeq{\status}{Attendee}{%
\DTLgadd{\attendee}{1}{\attendee}%
}{}%
}
\DTLsavelastrowcount{\totalconf}




\begin{cvhonors}
%\cvhonor{Carpentries Workshops}{Carpentries branded events and instructor training}{2017 -- Present}{\carpentries{}}
%\cvhonor{Long workshops}{}{2017 -- Present}{\longpres{}}
\cvhonor{Organizer}{}{2020 -- Present}{\organizer{}}
\cvhonor{Chair Positions}{}{2018 -- Present}{\chair{}}
\cvhonor{Program Committee Member/Reviewer}{}{2018 -- Present}{\reviewer{}}
\cvhonor{Community-Based Conference Attendee}{}{2013 -- Present}{\attendee{}}
\cvhonor{Total Participation}{}{2013 -- Present}{\totalconf{}}
\end{cvhonors}



\cvsubsection{Organizer}


\begin{cventries}
\DTLforeach[\equal{\generalRole}{Organizer}]{organizing}{%
%ProjectName	Date	Role	Location	Description1	Description2    
\projectname=ConferenceName,%
\date=Date,%
\role=Role,%
\generalRole=GeneralRole,%
\location=Location,%
\link=link}%
{%
 \cventry
    {} % subhead
    {\href{\link}{\projectname}} % head
    {\date} % green
    {\location}% % flushright grey
    {%
    }
}
\end{cventries}


\cvsubsection{Chair Positions}


\begin{cventries}
\DTLforeach[\equal{\generalRole}{Chair}]{organizing}{%
\projectname=ConferenceName,%
\date=Date,%
\role=Role,%
\generalRole=GeneralRole,%
\location=Location,%
\link=link}%
{%
 \cventry
    {\role} % subhead
    {\projectname} % head
    {\date} % green
    {\location} % flushright grey
    {%
    }
}
\end{cventries}

\vspace{10pt}


\cvsubsection{Program Committee Member}

\begin{cventries}
\DTLforeach[\equal{\role}{Program Committee Member}]{organizing}{%
\projectname=ConferenceName,%
\date=Date,%
\role=Role,%
\generalRole=GeneralRole,%
\location=Location,%
\link=link,
\type=type}%
{%
 \cventry
    {\type} % subhead
    {\projectname} % head
    {\date} % green
    {\location} % flushright grey
    {}%
}
\end{cventries}


\cvsubsection{Journal Reviewer}
\begin{cventries}
\DTLforeach[\equal{\role}{Journal Reviewer}]{organizing}{%
\projectname=ConferenceName,%
\date=Date,%
\role=Role,%
\generalRole=GeneralRole,%
\location=Location,%
\link=link}%
{%
    \begin{cvreviews}
    \cvreview
    {\projectname} % head
    {\date} % green
    \end{cvreviews}
}
\end{cventries}
\vspace{10pt}

\cvsubsection{ACL Rolling Review}
\begin{cventries}
\DTLforeach[\equal{\role}{ACL Rolling Review}]{organizing}{%
\projectname=ConferenceName,%
\date=Date,%
\role=Role,%
\generalRole=GeneralRole,%
\location=Location,%
\link=link}%
{%
    \begin{cvreviews}
      \cvreview
        {\generalRole}% role
        {\date}% Date(s)
    \end{cvreviews}
}
\end{cventries}
\vspace{10pt}


\cvsubsection{Community-Based Conference Attendance}
\begin{cventries}
\DTLforeach[\equal{\generalRole}{Attendee}]{organizing}{%
\projectname=ConferenceName,%
\date=Date,%
\role=Role,%
\generalRole=GeneralRole,%
\location=Location,%
\link=link}%
{%
    \begin{cvreviews}
    \cvreview
    {\projectname} % head
    {\date} % green
    \end{cvreviews}
}
\end{cventries}

\vspace{.5cm}
%-------------------------------------------------------------------------------
%	SECTION TITLE
%-------------------------------------------------------------------------------

{\color{black}\fontsize{12pt}{1em}\faIcon{chart-bar}} \cvsection{Professional Activities - Presentations}

\DTLloaddb[]{presentations}{cv/presentations.csv}
\cvsubsection{Summary}
%\def\invitedpres{0}
\def\campuspres{0}
\def\conferencepres{0}
\def\lecture{0}
\def\panel{0}
\def\totalpres{0}


\DTLforeach{presentations}{%
\status=Status}{%
\DTLifeq{\status}{Campus}{%
\DTLgadd{\campuspres}{1}{\campuspres}%
}{}%
\DTLifeq{\status}{Conference}{%
\DTLgadd{\conferencepres}{1}{\conferencepres}%
}{}%
\DTLifeq{\status}{Guest Lecture}{%
\DTLgadd{\lecture}{1}{\lecture}%
}{}%
\DTLifeq{\status}{Panel}{%
\DTLgadd{\panel}{1}{\panel}%
}{}%
}
\DTLsavelastrowcount{\totalpres}



\begin{cvhonors}
\cvhonor{Special Conference Presentations}{}{}{\conferencepres{}}
\cvhonor{On-Campus Invited Speaker Presentations}{}{}{\campuspres{}}
% \cvhonor{Guest Lectures}{}{}{\lecture{}}
\cvhonor{Panels}{}{}{\panel{}}
\cvhonor{Total Presentations}{}{}{\totalpres{}}

% \cvhonor{On-Campus Invited Speaker Presentations}{}{2013 -- Present}{\campuspres{}}
% \cvhonor{Special Conference Presentations}{}{2019 -- Present}{\conferencepres{}}
% \cvhonor{Guest Lectures}{}{2021 -- Present}{\lecture{}}
% \cvhonor{Total Presentations}{}{2013 -- Present}{\totalpres{}}
\end{cvhonors}

\vspace{0.7cm}


\newcommand*{\uniquetalks}{}

\cvsubsection{Special Conference Presentations}

\DTLforeach[\equal{\acadstatus}{Conference}]{presentations}{%
\projectname=ProjectName,%
\date=Date,%
\role=Role,%
\location=Location,%
\address=address,%
\acadstatus=Status,%
\link=link,%
\note=note}%
{%
    \begin{tabular*}{\textwidth}{@{\extracolsep{\fill}} L{\textwidth - 5cm} R{5cm}}
     \nohyphens{\entrytitlestyle{``\href{\link}{\projectname}''}}
    \end{tabular*}
    \newline
     \presentry% Usage: \cvsubentry{<position>}{<title>}{<date>}{<description>}
        {\date} % green
        {\location} % head
        {\address}% % flushright grey
        {}%
        \DTLifnullorempty
        {\note}
        {}
        {
        \begin{cvtalknotes}
            {\item {\bf \note}}
        \end{cvtalknotes}
        \vspace{5mm}
        }
}



\DTLforeach[\equal{\acadstatus}{Campus}]{presentations}{%
\prestitle=ProjectName,%
\date=Date,%
\role=Role,%
\location=Location,%
\address=address,%
\acadstatus=Status,%
\link=link,%
\note=note}%
{%
  \expandafter\DTLifinlist\expandafter{\prestitle}{\uniquetalks}%
  {}% do nothing, already in list
  {% add to list
    \ifdefempty{\uniquetalks}%
    {\let\uniquetalks\prestitle}% first element of list
    {% append to list
      \eappto\uniquetalks{,\prestitle}%
    }%
  }%
}

\vspace{0.6cm}


\cvsubsection{Invited Speaker Presentations}


\psforeach{\prestitle}{\uniquetalks}{% for each unique talk title
    \begin{tabular*}{\textwidth}{@{\extracolsep{\fill}} L{\textwidth - 5cm} R{5cm}}
    \nohyphens{\entrytitlestyle{``\prestitle''}}%
    \end{tabular*}%
    \newline
    \DTLforeach[\equal{\projectname}{\prestitle}\AND\equal{\acadstatus}{Campus}]{presentations}{% for each presentation, if it matches the title AND it is a campus talk
    \projectname=ProjectName,%
    \date=Date,%
    \role=Role,%
    \location=Location,%
    \address=address,%
    \acadstatus=Status,%
    \link=link,%
    \note=note}%
    {%
     \noindent\presentry% Usage: \cvsubentry{<position>}{<title>}{<date>}{<description>}
        {\date} % green
        {\href{\link}{\location}} % head
        {\address}% % flushright grey
        {}%\role}%
    }
}


\cvsubsection{Panelist}
\vspace{-.3em}
\begin{cventries}
\DTLforeach[\equal{\acadstatus}{Panel}]{presentations}{%
    \projectname=ProjectName,%
    \date=Date,%
    \role=Role,%
    \location=Location,%
    \address=address,%
    \acadstatus=Status,%
    \link=link,%
    \note=note}%
{%
 \cventry
    {\location} % subhead
    {\projectname} % head
    {\date} % green
    {\address} % flushright grey
    {%
    }
}
\end{cventries}


% \vspace{-.1cm}

% \cvsubsection{Guest Lectures}

% \DTLforeach[\equal{\acadstatus}{Guest Lecture}]{presentations}{%
% \projectname=ProjectName,%
% \date=Date,%
% \role=Role,%
% \location=Location,%
% \address=address,%
% \acadstatus=Status,%
% \link=link,%
% \note=note}%
% {%
%     % \begin{tabular*}{\textwidth}{@{\extracolsep{\fill}} L{\textwidth - 5cm} R{5cm}}
%     %  \nohyphens{\entrytitlestyle{``\href{\link}{\projectname}''}}
%     % \end{tabular*}
%     % \newline
    
%     \presentry% Usage: \cvsubentry{<position>}{<title>}{<date>}{<description>}
%         {\date} % green
%         {\location} % head
%         {\address}% % flushright grey
%         {}%
% }










% \DTLforeach[\equal{\acadstatus}{Guest Lecture}]{presentations}{%
% \projectname=ProjectName,%
% \date=Date,%
% \role=Role,%
% \location=Location,%
% \address=address,%
% \acadstatus=Status,%
% \link=link,%
% \note=note}%
% {%
%     \begin{tabular*}{\textwidth}{@{\extracolsep{\fill}} L{\textwidth - 5cm} R{5cm}}
%      \nohyphens{\entrytitlestyle{``\href{\link}{\projectname}''}}
%     \end{tabular*}
%     \newline
%     \presentry% Usage: \cvsubentry{<position>}{<title>}{<date>}{<description>}
%         {\date} % green
%         {\location} % head
%         {\address}% % flushright grey
%         {}%
% }



\vspace{.5cm}
%-------------------------------------------------------------------------------
%	SECTION TITLE
%-------------------------------------------------------------------------------
{\color{black}\fontsize{12pt}{1em}\faIcon{hands-helping}} \cvsection{Mentorship}

\DTLloaddb[]{mentorship}{cv/mentoring.csv}
%-------------------------------------------------------------------------------
%	CONTENT
%-------------------------------------------------------------------------------

\cvsubsection{Outreach}

\begin{cventries}
 \cventry
    {Out in Tech U’s Mentorship Program} % subhead
    {Mark McGovern}% head
    {Summer 2022}% green
    {}% flushright grey
    {}%
\end{cventries}
\hspace{2cm}

\cvsubsection{University of Pennsylvania}

\begin{cventries}
\DTLforeach[\equal{\location}{University of Pennsylvania}]{mentorship}{%
\person=person,%
\date=Date,%
\location=location,%
\then=Then,%
\now=Now}%
{%
 \cventry
    {\then \DTLifnullorempty{\now}{}{ $\rightarrow$ \now}} % subhead
    {\person}% head
    {\date}% green
    {}% flushright grey
    {}%
}%
\end{cventries}

\hspace{2cm}

\cvsubsection{Georgia Institute of Technology}

\begin{cventries}
\DTLforeach[\equal{\location}{Georgia Institute of Technology}]{mentorship}{%
\person=person,%
\date=Date,%
\location=location,%
\then=Then,%
\now=Now}%
{%
 \cventry
    {\then \DTLifnullorempty{\now}{}{ $\rightarrow$ \now}} % subhead
    {\person}% head
    {\date}% green
    {}% flushright grey
    {}%
}%
\end{cventries}

\vspace{.5cm}
%-------------------------------------------------------------------------------
%	SECTION TITLE
%-------------------------------------------------------------------------------
{\color{black}\fontsize{12pt}{1em}\faIcon{newspaper}} \cvsection{Media Coverage}


%-------------------------------------------------------------------------------
%	SUBSECTION TITLE
%-------------------------------------------------------------------------------



%-------------------------------------------------------------------------------
%	CONTENT
%-------------------------------------------------------------------------------
%\begin{cvmediaenv}


%---------------------------------------------------------
  \cvmedia
    {``\href{https://www.theregister.com/2023/08/19/chatgpt_dnd_dm/}{Hallucinating ChatGPT finds a role playing Dungeons \& Dragons}''} % title & link
    {Thomas Claburn} % author
    {The Register} % Location
    {Aug 19, 2023} % Date(s)

%---------------------------------------------------------
  \cvmedia
    {``\href{https://magazine.seas.upenn.edu/fall-2022/next-gen-innovators/}{Next-Gen Innovators: Penn Engineering Postdoctoral Fellows Lead the Way on Groundbreaking Research}''} % title & link
    {Amy Biemiller} % author
    {Penn Engineering Magazine} % Location
    {Fall 2022} % Date(s)

%---------------------------------------------------------
  \cvmedia
    {``\href{https://magazine.tank.tv/issue-88/talk/lara-martin}{Lara Martin: ``How can I get a system to tell a story about anything I want?''}''} % title & link
    {Masoud Golsorkhi} % author
    {TANK Magazine, Issue 88 (Narrative)} % Location
    {Autumn 2021} % Date(s)

%---------------------------------------------------------
  \cvmedia
    {``\href{https://www.sciencefocus.com/future-technology/lara-martin-on-teaching-ai-to-tell-stories/}{Lara Martin on teaching AI to tell stories}''} % title & link
    {Amy Barrett} % author
    {BBC Science Focus Podcast} % Location
    {Mar 15, 2021} % Date(s)

%---------------------------------------------------------
  \cvmedia
    {``\href{https://www.sciencefocus.com/magazine/dark-stars/}{Alexa, tell me a story}''} % Award
    {Amy Barrett} % author
    {BBC Science Focus Magazine} % Location
    {Feb 17, 2021} % Date(s)

%---------------------------------------------------------
  \cvmedia
    {``\href{https://mlatgt.blog/2020/04/16/meet-mlgt-lara-j-martin-trains-ai-agents-to-become-storytellers/}{Meet ML$@$GT: Lara J. Martin Trains AI Agents to Become Storytellers}''} % title & link
    {Allie McFadden} % author
    {Georgia Tech Machine Learning, Memos from ML\@GT} % Location
    {Apr 16, 2020} % Date(s)

%---------------------------------------------------------
  \cvmedia
    {``\href{https://www.wired.com/story/forget-chess-real-challenge-teaching-ai-play-dandd/}{Forget Chess---the Real Challenge Is Teaching AI to Play D\&D}''} % title & link
    {Will Knight} % author
    {Wired} % Location
    {Feb 28, 2020} % Date(s)

%---------------------------------------------------------
  \cvmedia
    {``\href{https://www.ic.gatech.edu/news/632082/changing-conversation-georgia-tech-researchers-provide-new-approach-automated-story}{Changing the Conversation: Georgia Tech Researchers Provide New Approach to Automated Story Generation}''} % title & link
    {David Mitchell} % author
    {Georgia Tech School of Interactive Computing} % Location
    {Feb 4, 2020} % Date(s)

%---------------------------------------------------------
  \cvmedia
    {``\href{https://zeszytymaryny.pl/naukowo/sztuczna-inteligencja-jako-pisarz-generowanie-fabuly/}{Sztuczna inteligencja jako pisarz: Generowanie fabuły}'' \newline (Translation from Polish: Artificial Intelligence as a Writer: Story Generation)} % title & link
    {Patrycja Świeczkowska} % author
    {Zeszyty Maryny (Blog)} % Location
    {Oct 4, 2019} % Date(s)

%---------------------------------------------------------

  \cvmedia
    {``\href{https://spectrum.ieee.org/tech-talk/robotics/artificial-intelligence/pictionary-playing-ai-sketches-the-future-of-human-machine-collaborations}{Pictionary-Playing AI Sketches the Future of Human-Machine Collaborations}'' (mention)} % title & link
    {Eliza Strickland} % author
    {IEEE Spectrum} % Location
    {Feb 6, 2019} % Date(s)

%---------------------------------------------------------

  \cvmedia
    {``\href{https://gvu.gatech.edu/georgia-tech-aaai2018}{Georgia Tech Artificial Intelligence Research Includes Collaborative Approaches with Humans, Automating Content, and More}''} % title & link
    {Joshua Preston} % author
    {Georgia Tech GVU Center} % Location
    {Feb 2, 2018} % Date(s)

%---------------------------------------------------------

%   \cvmedia
%     {``\href{https://medium.com/@mark_riedl/improvisational-storytelling-in-open-world-2359d7f649cf}{Improvisational Computational Storytelling in Open Worlds}''} % Award
%     {Mark Riedl} % Event
%     {Medium} % Location
%     {Jul 24, 2017} % Date(s)

%---------------------------------------------------------
%\end{cvmediaenv}



% {\color{black}\fontsize{12pt}{1em}\faIcon{users}} \cvsection{Institutional Service}

\begin{cventries}
\cventry
    {Reviewer}
    {President's Undergraduate Research Awards (PURA)}
    {Summer 2019}
    {Georgia Institute of Technology}
    {}    
\cventry
    {Volunteer}
    {School of Interactive Computing’s Prospective Student Visit Week}
    {Spring '16, '17, '18}
    {Georgia Institute of Technology}
    {}
\cventry
    {Member}
    {School of Interactive Computing Faculty Hiring Committee}
    {Fall 2018}
    {Georgia Institute of Technology}
    {}
\cventry
    {Member}
    {Graduate Student Council}
    {Fall 2018-Spring 2019}
    {Georgia Institute of Technology}
    {}    
\cventry
    {Website Manager}
    {Human-Centered Computing Website}
    {Fall 2017-Spring 2019}
    {Georgia Institute of Technology}
    {}    
\cventry
    {Coordinator}
    {School of Interactive Computing's Prospective Student Visit Week}
    {Spring 2016}
    {Georgia Institute of Technology}
    {}
\end{cventries}



%\input{cv/skills.tex}
% %-------------------------------------------------------------------------------
%	SECTION TITLE
%-------------------------------------------------------------------------------
{\color{black}\fontsize{12pt}{1em}\faIcon{building}} \cvsection{Research Experience}


%-------------------------------------------------------------------------------
%	CONTENT
%-------------------------------------------------------------------------------
\begin{cventries}


%---------------------------------------------------------
  \cventry
    {University of Maryland, Baltimore County -- Computer Science and Electrical Engineering (CSEE) Department} % Organization
    {Assistant Professor} % Job title
    {Aug 2023 – Present} % Date(s)
    {Baltimore, MD} % Location
    {
      % \begin{cvitems} % Description(s) of tasks/responsibilities
      % \end{cvitems}
    }

%---------------------------------------------------------
  \cventry
    {University of Pennsylvania -- Computer and Information Science} % Organization
    {\href{https://cifellows2020.org/}{Computing Innovation Fellow (CIFellow)} Postdoctoral Researcher} % Job title
    {Jan 2021 – Aug 2023} % Date(s)
    {Philadelphia, PA} % Location
    {
      % \begin{cvitems} % Description(s) of tasks/responsibilities
      % \item {Identifying the story understanding capabilities of large language models (LLMs).}
      % \item {Developing a working AAC prototype given feedback from users.}
      % \item {Conducted semi-structured interviews with autistic adult users of augmentative and alternative communication (AAC).}
      % \end{cvitems}
    }
%---------------------------------------------------------
  \cventry
    {Georgia Institute of Technology -- School of Interactive Computing} % Organization
    {Graduate Research Assistant} % Job title
    {Aug 2015 – Dec 2020} % Date(s)
    {Atlanta, GA} % Location
    {
      % \begin{cvitems} % Description(s) of tasks/responsibilities
      %   \item {Created a complex end-to-end automated story generation pipeline.}
      % \end{cvitems}
    }

%---------------------------------------------------------
  \cventry
    {Amazon.com Inc. -- Alexa Smart Home Machine Learning} % Organization
    {Applied Scientist Intern} % Job title
    {May 2017 – Aug 2017} % Date(s)
    {Seattle, WA} % Location
    {
      % \begin{cvitems} % Description(s) of tasks/responsibilities
      %   \item {Identified potential research questions within Alexa Smart Home.}
      %   \item {Developed a system for identifying commands with an assumed context.}
      % \end{cvitems}
    }

%---------------------------------------------------------
  \cventry
   {Carnegie Mellon University -- Language Technologies Institute} % Organization
    {Graduate Research Assistant} % Job title
    {Sept 2013 – Aug 2015} % Date(s)
    {Pittsburgh, PA} % Location
    {
      % \begin{cvitems} % Description(s) of tasks/responsibilities
      %   \item {Created a zero-resource speech-to-speech translation system for the University of Pittsburgh Medical Center.}
      %   \item {Performed emotion recognition in noisy speech for event detection.}
      % \end{cvitems}
    }
%---------------------------------------------------------
  \cventry
   {University of Southern California -- Institute for Creative Technologies} % Organization
    {Intern} % Job title
    {May 2011 – Aug 2011} % Date(s)
    {Playa Vista, CA} % Location
    {
      % \begin{cvitems} % Description(s) of tasks/responsibilities
      %   \item {Wrote a chatbot for the Virtual Patient Project using Bruce Wilcox’s language Chatscript.}
      %   \item {Developed an authoring tool for the Chatscript language using Java.}
      %   \item {Designed and ran experiments comparing my Chatscript system to the project’s current chat system.}
      % \end{cvitems}
    }

%---------------------------------------------------------
\end{cventries}

% \input{cv/extracurricular.tex}
% \input{cv/honors.tex}
% \input{cv/presentation.tex}
% \input{cv/writing.tex}




\vfill
{\tiny
CV produced using \LaTeX. Modified version of code: \href{https://github.com/Denubis/Academic-LaTeX-CV-BibTeX-CSV}{\tt github.com/Denubis/Academic-LaTeX-CV-BibTeX-CSV}.}

%-------------------------------------------------------------------------------
\end{document}
